\section{Results: data-snooping tests}
\label{sec:snooping}

\subsection{The studentized SPA test is oversized; the Reality Check is not}
Hansen's studentized test for superior predictive ability, as Section~\ref{sec:design} defines it, is oversized at two years of daily data and twenty strategies (Table~\ref{tab:size-joint}). Across the nine families and both governed runs its consistent $p$-value rejected \nz{both.size_range.spa_consistent}\% of skill-less searches at a nominal 5\%: \nz{v1.size.spa_consistent.iid_normal}\% under i.i.d.\ normal returns, \nz{v1.size.spa_consistent.skew_negative}\% under negative skew and \nz{v1.size.spa_consistent.ar1}\% under autocorrelation in v1. It exceeded the pre-registered 6\% bar in \nz{v1.fails6.spa_consistent} of the nine families in v1. The upper $p$-value behaves almost identically, and StepM, whose first step uses the same studentized statistic and recentring, matched the consistent SPA test to within \nz{both.stepm_spa_maxdiff} percentage points in every family of both runs.

White's Reality Check, which uses the same bootstrap but does not studentize, rejected \nz{both.size_range.reality_check}\% across the same cells, and \nz{both.size_range_nonar.reality_check}\% outside the AR(1) family. Among the v1 validators it is the only one that is calibrated, or close to it, under negative skew, positive skew, fat tails, volatility clustering and correlated strategies at once.

This was not the comparison v1 was designed to make. Its pre-registration listed the two SPA $p$-values as candidates for the author's software and only measured the Reality Check; the pre-registered rule therefore chose the consistent SPA test, with failures in most families. Rather than switch to the Reality Check on that evidence, a second pre-registration (v1b) added it as the first candidate and repeated the run on fresh seeds. v1b confirmed the pattern (Table~\ref{tab:size-joint}); its pre-registered rule chose the Reality Check, which missed the bar only in \nz{v1b.headline.joint.fails}.

The pattern does not depend on the implementation. In v2, an independent Python implementation written from the published definitions reproduced v1b's sizes of the consistent SPA test and the Reality Check in all nine families to within three Monte Carlo standard errors (prediction P1 \nz{v2.P1.held}; the largest difference was \nz{v2.P1.maxdiff} points).

\subsection{Why: a variance estimate the bootstrap holds fixed}
The studentized statistic divides each mean by $\hat\omega_k$, an estimate of its long-run standard deviation made once from the sample, and the resampled statistics divide by the same $\hat\omega_k$. The bootstrap therefore reproduces the sampling variation of the means but not that of the estimate in the denominator. \citet{hall1991guidelines} give the general rule this breaks: a resampled statistic should be built the same way as the observed one. The Politis--Romano estimator makes the omission costly, because at the bootstrap's block length it adds weighted sample autocovariances to the variance. Under i.i.d.\ normal returns its standard deviation is, to first order, \nz{approx.cv_w.b8}\% of $\omega$ at block length 8 and \nz{approx.cv_w.b22}\% at 22, against \nz{approx.cv_w.b1}\% for the sample standard deviation. A studentized mean then behaves roughly like Student $t$ with about \nz{approx.nu.b8} (block 8) or \nz{approx.nu.b22} (block 22) degrees of freedom, while the bootstrap's critical value for the best of \nz{trials} strategies stays near the normal value of \nz{approx.z_star}. This back-of-the-envelope account predicts sizes of \nz{approx.size.b8}\% and \nz{approx.size.b22}\%; it gets the ordering right and accounts for \nz{approx.share.b8}\% and \nz{approx.share.b22}\% of the excess measured in experiment A, the rest coming from what it ignores (the skewness of the variance estimate's own distribution, and the bootstrap's noise).

Experiment A tests the account directly (Table~\ref{tab:mechanism}). Studentizing by the sample standard deviation instead, with the bootstrap unchanged, removes most of the excess where returns are symmetric and serially uncorrelated: \nz{v2.A.size.spa_c_sd.iid_normal}\% under i.i.d.\ normal returns, \nz{v2.A.size.spa_c_sd.student_t4}\% under $t_4$, \nz{v2.A.size.spa_c_sd.garch}\% under GARCH. It does not remove the excess under negative skew (\nz{v2.A.size.spa_c_sd.skew_negative}\%), where a second effect appears: a high sample mean tends to come with a low sample standard deviation, so the studentized statistic has a heavier upper tail than any fixed denominator can reproduce. A bootstrap-$t$ that re-studentizes every resample by its own standard deviation reproduces both effects and is calibrated under negative skew (\nz{v2.A.size.spa_c_boot_t.skew_negative}\%); outside the AR(1) family its size was \nz{v2.A.size_range_nonar.spa_c_boot_t}\% (predictions P3 and P4 \nz{v2.P3.held} and \nz{v2.P4.held}). With the long block arch uses by default ($\lfloor\sqrt{T}\rfloor = 22$), the fixed-variance studentized test rejected \nz{v2.A.size.spa_c_b22.iid_normal}\% of i.i.d.\ normal searches and \nz{v2.A.size.spa_c_b22.skew_negative}\% under negative skew (P6 \nz{v2.P6.held}).

One pre-registered prediction about the mechanism failed. P2 expected the sample-standard-deviation version to stay at or below 6\% in \nz{v2.P2.family_count} families; it did in \nz{v2.P2.pass_count}, but rejected \nz{v2.A.size.spa_c_sd.block_cluster}\% in the block-cluster family. In that family every joint test was slightly liberal in v2, including the unstudentized ones (Reality Check \nz{v2.A.size.rc.block_cluster}\%, bootstrap-$t$ \nz{v2.A.size.spa_c_boot_t.block_cluster}\%), as were all the joint tests with only five i.i.d.\ normal strategies in experiment B (Reality Check \nz{v2.B.size.rc.iid_normal.n504_k5}\%). A search of twenty strategies in four tight clusters has few effectively independent strategies, and the bootstrap's own noise, which matters less when a maximum over many strategies dominates the statistic, is then visible.

\subsection{How the excess depends on the sample and the search}
The excess is a finite-sample effect: it shrinks as the sample grows and grows with the number of strategies (Table~\ref{tab:grid}; prediction P7 \nz{v2.P7.held}). Under i.i.d.\ normal returns the fixed-variance test rejected \nz{v2.B.size.spa_c.iid_normal.n252_k20}\% of skill-less searches with one year of daily data, \nz{v2.B.size.spa_c.iid_normal.n504_k20}\% with two and \nz{v2.B.size.spa_c.iid_normal.n2520_k20}\% with ten; with two years of data it rejected \nz{v2.B.size.spa_c.iid_normal.n504_k5}\% of searches of five strategies and \nz{v2.B.size.spa_c.iid_normal.n504_k100}\% of searches of a hundred. Under negative skew the same figures were \nz{v2.B.size.spa_c.skew_negative.n252_k20}\%, \nz{v2.B.size.spa_c.skew_negative.n504_k20}\%, \nz{v2.B.size.spa_c.skew_negative.n2520_k20}\%, \nz{v2.B.size.spa_c.skew_negative.n504_k5}\% and \nz{v2.B.size.spa_c.skew_negative.n504_k100}\%. The bootstrap-$t$ and the Reality Check stayed near 5\% across the grid (\nz{v2.B.range.spa_c_boot_t}\% and \nz{v2.B.range.rc}\% over all twelve cells). \citet{hsu2010stepwise} report family-wise error rates of the studentized stepwise SPA test controlled at the nominal level in simulations with 1,000 periods and between 90 and 9,000 strategies; their designs differ from ours in sample length, block length and the mix of neutral and poor strategies, and we have not reproduced them. Distinct from this finite-variance effect, \citet{frederiksen2026heavy} show that when loss differentials have infinite variance, conventional predictive-ability tests can be severely oversized even asymptotically.

\subsection{What this means for software}
The open-source Python package arch \citep{sheppard2025arch} does not studentize its SPA statistic. In version \nz{v2.D.arch} the \texttt{studentize} option changes neither the statistic nor the resampled values; the variance estimate enters only the consistent recentring threshold. On \nz{v2.D.total_pairs} search-and-block-length pairs in experiment D, arch's $p$-values were identical with \texttt{studentize=True} and \texttt{studentize=False} on \emph{every} pair (\nz{v2.D.total_flag_changes} differed), and they agreed with this paper's unstudentized consistent test to within Monte Carlo error (mean absolute difference \nz{v2.D.meanabs_range}; Table~\ref{tab:arch}). Arch's ``SPA'' test is thus the Reality Check with Hansen's consistent recentring, and in the Null Zoo it stays close to its level outside the AR(1) family (\nz{v2.A.size_range_nonar.spa_c_unstud}\% for our implementation of the same computation). A pull request opened on 27 September 2026 proposes to make \texttt{studentize=True} divide by the fixed standard error, at arch's default block length \citep{arch871}. Apart from the floor at zero, which rarely matters under these null hypotheses, that is the computation measured above at \nz{v2.A.size.spa_c_b22.iid_normal}\% and \nz{v2.A.size.spa_c_b22.skew_negative}\%; we implemented it independently from the formula and did not run the pull request's code. Studentizing inside the bootstrap would avoid the excess.

A second, smaller issue is a matter of convention. With the statistic floored at zero, $T^{\mathrm{SPA}}=0$ whenever no strategy beats the benchmark in sample. If every strategy also falls below the consistent threshold, the recentred resampled statistics are mostly zero, and a $p$-value computed with a strict inequality, $\Pr^*(T^*>T)$, counts only the few resamples above zero: the test rejects although no strategy did better than the benchmark. This caused the second failed prediction. P9 expected the bootstrap-$t$ version to keep its size in a design whose other strategies were clearly poor; it rejected \nz{v2.C.size.spa_c_boot_t.poor_alternatives}\%, as did every floored variant, while arch's unfloored computation rejected \nz{v2.C.size.spa_c_unstud.poor_alternatives}\%. An exploratory check that was not pre-registered (fresh seed \nz{explore.seed}, \nz{explore.searches} searches) confirmed the cause: the observed statistic was zero in \nz{explore.share_T_zero}\% of searches and every strategy fell below the threshold in \nz{explore.share_all_below}\%, which contributed \nz{explore.size_T_zero.spa_c_boot_t_strict} points of the bootstrap-$t$ test's \nz{explore.size.spa_c_boot_t_strict}\%. Counting ties, $\Pr^*(T^*\ge T)$, which returns $p=1$ whenever $T=0$, brought it to \nz{explore.size.spa_c_boot_t_ties_count}\%. The implementation behind v1, the author's, also uses the strict inequality; in v1's null cells, where every strategy has mean zero, a search in which every strategy falls below the threshold is too rare to matter.

\subsection{Where studentizing pays off}
Studentizing exists for a reason, and experiment C shows it (Table~\ref{tab:fair}). When the strategies' volatilities range from 0.5 to 2 and the skilled strategy is the least volatile, the largest mean is the skilled strategy's in only \nz{v2.C.best_mean.heterogeneous_volatility}\% of searches, while the largest $t$-statistic is in \nz{v2.C.best_t.heterogeneous_volatility}\%. The Reality Check, which ranks by mean, then rejects in \nz{v2.C.power.rc.heterogeneous_volatility}\% of searches with a Sharpe ratio of 2: no more than under the null. Arch's unstudentized computation does no better (\nz{v2.C.power.spa_c_unstud.heterogeneous_volatility}\%). The studentized tests find the strategy about half the time (bootstrap-$t$ \nz{v2.C.power.spa_c_boot_t.heterogeneous_volatility}\%, size-adjusted \nz{v2.C.adj.spa_c_boot_t.heterogeneous_volatility}\%), and the fixed-variance version is again oversized (\nz{v2.C.size.spa_c.heterogeneous_volatility}\%). When the other strategies are clearly poor, the Reality Check, which recentres every strategy at its mean, is very conservative (size \nz{v2.C.size.rc.poor_alternatives}\%, power \nz{v2.C.power.rc.poor_alternatives}\%), and Hansen's consistent recentring raises power to \nz{v2.C.power.spa_c_unstud.poor_alternatives}\% at a size of \nz{v2.C.size.spa_c_unstud.poor_alternatives}\% (unstudentized) (P8 \nz{v2.P8.held}). The lesson is not to give up studentizing or recentring, which is why Hansen proposed them, but to studentize inside the bootstrap and to count ties.

\subsection{Power}
Where the strategies' volatilities are equal, as in the nine families, studentizing buys no power: at a Sharpe ratio of 2 the bootstrap-$t$ version and the Reality Check had similar size-adjusted power under i.i.d.\ normal returns (\nz{v2.A.adj.spa_c_boot_t.iid_normal}\% and \nz{v2.A.adj.rc.iid_normal}\%; Table~\ref{tab:mechanism-power}). Skew pulls them apart in opposite directions: under negative skew the Reality Check was ahead (\nz{v2.A.adj.rc.skew_negative}\% against \nz{v2.A.adj.spa_c_boot_t.skew_negative}\%), under positive skew the bootstrap-$t$ (\nz{v2.A.adj.spa_c_boot_t.skew_positive}\% against \nz{v2.A.adj.rc.skew_positive}\%). The fixed-variance test's apparent power advantage (raw \nz{v1.power2.spa_consistent.iid_normal}\% in v1, against \nz{v1.power2.reality_check.iid_normal}\% for the Reality Check) disappears once its excess size is removed (\nz{v1.adj2.spa_consistent.iid_normal}\% against \nz{v1.adj2.reality_check.iid_normal}\%).

\subsection{The pre-registered predictions}
Of the \nz{v2.predictions.count} predictions registered before the v2 run, \nz{v2.predictions.held_count} held. Table~\ref{tab:predictions} lists all of them with their outcomes. The two that failed, P2 and P9, are reported above with their explanations; the explanation of P9 comes from an exploratory check made after the run and should be read as such.
